\chapter{Conclusão e Trabalhos Futuros}
\label{chap:conclusao}

\section{Síntese dos Resultados}

\orientacao{O que o projeto desenvolveu, organizado por pacote, em lista de resultados concretos. Cada item nomeia um artefato ou um achado, com o número que o caracteriza (quantidade de imagens, número de espécies, arquitetura usada). Feche com uma frase sobre a hipótese central: ela se confirmou, se confirmou em parte ou não se confirmou.

Declare também o nível de maturidade tecnológica atingido, na forma ``o projeto atinge o nível X, comprovado por Y'', com escala, fidelidade e ambiente, e o mesmo valor lançado no SRInfo.}

\begin{itemize}
    \item \ph{Pacote 1: artefato ou achado, com o número que o caracteriza.}
    \item \ph{Pacote 2: artefato ou achado, com o número que o caracteriza.}
\end{itemize}

\ph{A hipótese central do projeto se confirmou / se confirmou em parte / não se confirmou, porque ...} O projeto atinge o nível TRL~\ph{X}, comprovado por \ph{evidência}~\cite{abnt16290}.

\section{Limitações}

\orientacao{Uma limitação por parágrafo, com título em negrito e com causa e efeito. A primeira deve tocar o objetivo central do trabalho, quando houver limitação ali: registrá-la é o que dá autoridade ao documento, e o parceiro já a conhece.}

\paragraph{\ph{Nome da limitação}.} \ph{Causa técnica e efeito sobre o resultado.}

\paragraph{\ph{Nome da limitação}.} \ph{Causa técnica e efeito sobre o resultado.}

\section{Trabalhos Futuros}

\orientacao{Uma direção por parágrafo, com título em negrito, ligada a uma limitação ou a uma oportunidade identificada. Escreva como sugestão técnica, não como compromisso nem como proposta comercial. Se há intenção de um projeto seguinte, isso se conversa em reunião e entra em proposta, não em parágrafo de relatório.}

\paragraph{\ph{Direção}.} \ph{Sugestão técnica, ligada a uma limitação ou a uma oportunidade.}

\section{Considerações Finais}

\orientacao{Um parágrafo. Diz o que ficou demonstrado, o que o parceiro leva e o que o conjunto (conhecimento, dados, resultados experimentais) habilita. Sem superlativo e sem promessa de produto.}

\ph{Um parágrafo: o que ficou demonstrado, o que o parceiro leva e o que o conjunto habilita.}
